\section{Causal framework}\label{sec:context}


\subsection{Modelling of causal effects}

Let $T$ denote the non-negative survival time and $C$ denote the censoring time. 
We observe the follow-up time $Y = \min(T, C)$ along with the censoring indicator $\delta \in \{0,1\}$, where $\delta = 1$ if the event is observed and $\delta = 0$ if the observation is right-censored.
Let $A \in \{0,1\}$ be a binary treatment indicator, where $A = 1$ corresponds to the treated group and $A = 0$ to the control group. 
Covariates are denoted by $X \in \mathbb{R}^p$, where $p$ is the number of observed pre-treatment variables.
The covariates may include potential confounders.
We do not distinguish between various types of covariates in our notation, such as confounders, prognostic factors, moderators, or predictors of treatment assignment.
Throughout, we use capital letters (e.g., $T$, $A$, $X$) to denote random variables, and lowercase letters (e.g., $t$, $a$, $x$) to denote their realisations. 
We consider a sample of $n$ independent observations. 

We assume an accelerated failure time (AFT) model \citep{buckley1979linear} for the potential survival times. 
The AFT model relates the logarithm of the survival time to a regression function.
We decompose the log-transformed outcome into a \textit{prognostic} component and a \textit{treatment effect} component \citep{hahn2020bayesian}:
\begin{equation}\label{model}
    \log T(a) = f(x, \hat{e}(x)) + a \cdot \tau(x) + \varepsilon,
\end{equation}
where $\varepsilon \sim \mathcal{N}(0, \sigma^2)$ denotes the error term. 
The function $f(x, \hat{e}(x))$ models the baseline prognosis and incorporates the estimated propensity score $\hat{e}(x)$.
The propensity score, $e(x) = \mathbb{P}(A = 1 \mid X = x)$, is the probability of treatment given covariates \citep{RosenbaumRubin1983}.
The function $\tau(x)$ directly models the treatment effect.
In Section~\ref{section:Model}, we describe a novel regression framework for $f$ and $\tau$ that builds on Bayesian regression trees and incorporates a horseshoe prior to handle the high-dimensional covariates.
We refer to model \eqref{model} with horseshoe prior on the step heights of $f$ and $\tau$ as the \textsl{Causal Horseshoe Forest}.


The AFT model offers several advantages in the context of causal inference with survival data. 
The AFT model provides a \textit{collapsible} estimand: the acceleration factor, defined as the ratio of expected survival times under treatment and control, remains invariant when conditioning on additional non-confounding covariates \citep{Robins1992AFT, Aalen2015CoxCausal, Crowther2023}.
This implies that $\tau(x)$ can be meaningfully averaged across covariate distributions to obtain the marginal causal effect without introducing bias due to non-collapsibility.
Collapsibility is a desirable property in causal inference, as it ensures interpretability and comparability of marginal and conditional effects. 
In contrast, commonly used hazard-based measures, such as the hazard ratio in the Cox model, are non-collapsible \citep{Hernan2010, Martinussen2013, Aalen2015CoxCausal}.
Furthermore, \cite{Hougaard1999} showed that AFT models are robust to omitted covariates, enhancing their suitability for observational studies.

We include the estimated propensity score in the outcome model to improve causal effect estimation in high-dimensional settings. 
This inclusion helps to control for confounding and mitigates regularisation-induced confounding, where aggressive shrinkage of confounders can bias treatment effect estimates \citep{hahn2020bayesian, Ray2020SemiparametricBayes}. 
The propensity score acts as a balancing score: given the propensity score, covariates and treatment assignment become independent \citep{RosenbaumRubin1983}. 
We incorporate the \textit{estimated} propensity score directly in the prognostic regression function $f(x, \hat{e}(x))$ to avoid feedback between the outcome and treatment models.
Although this approach is not dogmatically Bayesian and does not propagate uncertainty from the propensity score estimation \citep{Zigler2013ModelFeedback}, it improves performance in finite samples \citep{hahn2020bayesian}. 
We estimate the propensity score using a tree-based model specifically designed for high-dimensional covariates and binary treatment assignment.



\subsection{Estimands and identification assumptions}

We define causal estimands within the potential outcomes framework \citep{Rubin1974, Rubin1978}. 
For each individual, let $T(1)$ and $T(0)$ represent the survival times that would be observed under treatment and control, respectively. 
Likewise, $C(1)$ and $C(0)$ denote the potential censoring times corresponding to each level of treatment.
We focus on the conditional average treatment effect (CATE), defined as the difference in expected log survival times under treatment and control, conditional on observed covariates:
\begin{equation}
    \text{CATE}(x) := \E[\log T(1) - \log T(0) \mid X = x].
\end{equation}
The CATE captures heterogeneity in treatment effects across subgroups defined by covariate profiles $X = x$.  
We emphasise that the CATE is a population-level functional of the conditional potential outcome distributions.
It should be interpreted as a function of covariates $x$ and not as a realised individual-level causal contrast, since it depends only on the marginal conditional expectations $\E[\log T(a)\mid X=x]$ for $a\in\{0,1\}$ and does not involve the joint distribution of potential outcomes.


We rely on common assumptions to identify $\tau(x)$ from observational data \citep{ImbensRubin2015}.
The first is the \textit{Stable Unit Treatment Value Assumption} (SUTVA), which rules out interference across units and presumes each individual has well-defined potential outcomes under both treatment conditions.
This implies that the observed survival time can be expressed as $T = A T(1) + (1 - A) T(0)$.
We assume \textit{unconfoundedness} (see \eqref{unconfoundedness}) and \emph{positivity}.
Positivity requires that for all covariate profiles $x $ each treatment arm has a non-zero probability of assignment:
\begin{equation}
0 < e(x)  < 1.
\end{equation}
The assumptions above are sufficient to identify the causal estimands in terms of the survival times $T$ conditional on $(X,A)$. 
An additional assumption on the censoring mechanism is required to identify these quantities from the observed, right-censored data.
We assume \emph{independent censoring}, meaning that the censoring mechanism is conditionally independent of the survival time given treatment and covariates:
\begin{equation}
C(a) \ind T(a) \mid X = x, A=a \quad\text{for } a \in \{0,1\}.
\end{equation}
Under this assumption, the observed data $(Y,\delta)$ are sufficient to identify the conditional distribution of $T$ given $(X,A)$ from right-censored observations \citep{Andersen1993, Lagakos1979}.
Conditional on $(X,A)$, the censoring mechanism does not depend on the event time.
The observed events are representative of the underlying event-time distribution within each covariate--treatment stratum.
As a result, functionals of the conditional distribution of $T$, including $\E[\log T \mid X=x,A=a]$, can be identified from the censored data.


The causal effect is identifiable from the observed data under these assumptions \citep{Imbens2004Identification}, i.e.  \protect{$\E[\log T(a) \mid X = x] = \E[\log T \mid X = x, A = a]$} for each $a \in \{0,1\}$.
As a result, the CATE can be expressed in terms of observable quantities as:
\begin{equation}
\begin{split}
    \text{CATE}(x) &= \E[\log T \mid X = x, A = 1] - \E[\log T \mid X = x, A = 0] =\tau(x).
\end{split}
\end{equation}
The final equality follows from model assumption \eqref{model}. 
This identification enables a direct modelling approach of the conditional average treatment effect.